\documentclass[11pt,twocolumn]{article}
\usepackage[margin=0.85in,top=1in]{geometry}

\usepackage{amsmath,amssymb,amsthm}
\usepackage{graphicx}
\usepackage{microtype}
\usepackage{hyperref}
\usepackage{xcolor}
\usepackage{booktabs}
\usepackage{enumitem}
\usepackage{titlesec}
\titleformat{\section}{\normalfont\bfseries\large}{\thesection}{1em}{}
\titleformat{\subsection}{\normalfont\bfseries\normalsize}{\thesubsection}{1em}{}
\titlespacing*{\section}{0pt}{14pt}{4pt}
\titlespacing*{\subsection}{0pt}{10pt}{3pt}

\hypersetup{colorlinks=true,linkcolor=blue,citecolor=blue,urlcolor=blue}

\theoremstyle{plain}
\newtheorem{theorem}{Theorem}[section]
\newtheorem{definition}[theorem]{Definition}
\newtheorem{proposition}[theorem]{Proposition}
\theoremstyle{remark}
\newtheorem{remark}[theorem]{Remark}

% ── Shortcuts ──────────────────────────────────────────────────────────
\newcommand{\DAG}{\mathcal{G}}
\newcommand{\partialG}{\partial\DAG}
\newcommand{\Vvalid}{\mathcal{V}_{\mathrm{valid}}}
\newcommand{\lambdax}{\lambda(x)}
\newcommand{\DeltaS}{\Delta S(\rho\|\sigma_0)}

\begin{document}

\title{Relational Actualism:\\
Irreversible Events as the Primitive Basis of Physical Reality}

\author{Joshua F.~Sandeman\\
\small Independent Researcher, Salem, Oregon, USA\\
\small \href{mailto:sansuikyo@gmail.com}{sansuikyo@gmail.com}}

\date{March 2026}

% ── ABSTRACT ───────────────────────────────────────────────────────────
\begin{abstract}
Relational Actualism (RA) is a foundational framework in which
on-shell actualization events---irreversible physical interactions
that write permanent vertices into a growing causal Directed
Acyclic Graph (DAG)---are taken as the primitive constituents
of physical reality.
From this single ontological commitment, no new fields or free
parameters are required to derive: the observer-independent
criterion for wave function collapse ($\Delta S(\rho\|\sigma_0)>0$);
the structural zero of the cosmological constant; the density-dependent
Hubble expansion rate; the prohibition of causally silent dark matter
candidates; flat galactic rotation curves; the fine structure constant
$\alpha_{\rm EM}^{-1}=137$ and strong coupling $\alpha_s(m_Z)=
1/\sqrt{72}$; and the complete particle content of the Standard Model
from five integers.
We present the Engine of Becoming---the five-step Covariant Stochastic
Graph Rewriting System that governs the growth of the causal DAG---and
show that quantum mechanics and general relativity are not rival
fundamental theories but descriptions of different levels and limits
of the same algorithm.
Time in RA is not a background parameter but the process of graph growth
itself: the arrow of time is not explained by entropy gradients but
encoded in the irreversibility of every actualization event.
All framework claims with finite-dimensional analogs have been formally
verified in Lean~4/Mathlib.
We articulate the hard mathematical walls that remain open,
identify the collaboration targets that would close them, and assemble
the complete falsifiable prediction set of the programme.
\end{abstract}

\maketitle

% ── I. INTRODUCTION ─────────────────────────────────────────────────────
\section{Introduction: The Shared Error}

The irreconcilability of quantum mechanics (QM) and general relativity
(GR) is routinely presented as the central unsolved problem of
theoretical physics.
Every proposed resolution---string theory, loop quantum gravity,
causal dynamical triangulations, asymptotic safety---accepts both
frameworks as given and attempts to make them mutually consistent.
We propose instead that the difficulty lies upstream: both frameworks
share a foundational error, and identifying that error dissolves the
tension rather than resolving it.

\emph{The shared error is the treatment of time as a background parameter.}
In the Schr\"{o}dinger equation, $t$ is a continuous coordinate along
which the state vector evolves unitarily.
In the Einstein field equations, $t$ is a coordinate label on a
block Lorentzian manifold.
In both cases, nothing actually happens: the formalism describes a
static mathematical structure; genuine temporal becoming is represented
nowhere.

Relational Actualism begins from a different premise.
Time is not a dimension through which events are located.
Time is the process of the causal graph growing---the irreversible
accumulation of on-shell actualization events that write permanent
vertices into a growing Directed Acyclic Graph (DAG).
The arrow of time is not a consequence of entropy gradients or
initial conditions; it is encoded in the very definition of an
actualization event, which is irreversible by stipulation.

This is not merely an interpretive shift.
The companion papers of this programme~\cite{Sandeman2026P1,Sandeman2026P2,Sandeman2026P3}
have shown that this single ontological reorientation yields, without
additional parameters: the correct values of the two most precisely
measured coupling constants; an observer-independent criterion for
wave function collapse with three testable predictions; a structural
resolution of the cosmological constant problem; and a geometric mechanism
for the Hubble tension with a falsifiable numerical prediction.

The present paper assembles the framework in its most complete form.
Section~\ref{sec:primitive} states the primitive ontological commitment
and its immediate logical consequences.
Section~\ref{sec:engine} presents the Engine of Becoming: the
five-step algorithm governing the growth of the causal DAG.
Section~\ref{sec:limits} shows how QM and GR emerge as limiting
approximations.
Section~\ref{sec:complexity} develops the emergence of biological
complexity and the Causal Firewall.
Section~\ref{sec:time} develops the account of time, causality, and
the arrow.
Section~\ref{sec:philosophy} positions the framework against the
principal competing ontologies.
Section~\ref{sec:predictions} assembles all testable predictions
of the programme.
Section~\ref{sec:hardwalls} states the open mathematical problems
with precision.

% ── II. THE PRIMITIVE COMMITMENT ────────────────────────────────────────
\section{The Primitive Commitment}
\label{sec:primitive}

\subsection{Actualization events as primitive}

Relational Actualism posits:

\begin{definition}[Actualization event]
An \emph{actualization event} is a physical interaction that produces
an irreversible increase in quantum relative entropy
$\Delta S(\rho\|\sigma_0) > 0$ with respect to the Minkowski vacuum
state $\sigma_0$, generating a permanent, redundant causal record
in the asymptotic environment.
\end{definition}

This is an objective, observer-independent physical criterion.
The relative entropy $S(\rho\|\sigma_0) = \mathrm{Tr}[\rho(\ln\rho
- \ln\sigma_0)]$ is a property of the interaction, not of any
observer's knowledge.
The vacuum state $\sigma_0$ is the physical reference state of
the quantum field, not a subjective baseline.
In weakly coupled quantum field theory, the criterion coincides
with the on-shell condition $p^\mu p_\mu = m^2c^2$: a real excitation
propagating into the asymptotic field produces exactly the permanent
environmental record required~\cite{Sandeman2026P2}.

\subsection{Three kinds of property}

Properties in RA are of exactly three kinds:

\begin{enumerate}[label=(\roman*), noitemsep]
\item \emph{Invariant properties}: possessed by all actualized states
  regardless of context (mass, charge, spin quantum number).
\item \emph{Classical properties}: actualized deterministically through
  interaction, because the causal and logical context has reduced the
  possibilities to exactly one.
\item \emph{Intrinsically relational properties}: actualized only in
  specific interaction contexts, with genuine ontic randomness when
  more than one possibility remains viable.
\end{enumerate}

This taxonomy dissolves the measurement problem at its root: asking
``what is the value of an intrinsically relational property prior to
the interaction that actualizes it?'' is not a question about Nature's
hidden state but a category error.
Intrinsically relational properties are not hidden; they are not yet
actual.

\subsection{The causal DAG}

Let $\DAG = (V, E)$ be a growing Directed Acyclic Graph in which
each vertex $v \in V$ is an actualization event and each edge
$(u, v) \in E$ encodes ``$u$ is in the causal past of $v$.''
The frontier $\partialG$ is the set of actualized vertices with
no causal successors: the dynamical present.
The graph grows by the sequential addition of new vertices to
$\partialG$.
The continuum spacetime manifold is the dense-graph limit of $\DAG$,
recovered in the regime $\lambda \to \infty$ where $\lambda$ is
the actualization density (vertices per unit proper volume).

\begin{definition}[Local Ledger Condition]
For every vertex $v \in V$, the sum of incoming conserved charges
equals the sum of outgoing charges:
\begin{equation}
  \sum_{\substack{u:\,(u,v)\in E}} \mathbf{q}(u)
  \;=\;
  \sum_{\substack{w:\,(v,w)\in E}} \mathbf{q}(w).
  \label{eq:LLC}
\end{equation}
\end{definition}

The LLC is the discrete graph-theoretic primitive from which
Noether conservation laws emerge in the continuum limit.
It is formally verified in Lean~4/Mathlib (theorem
\texttt{local\_ledger\_condition} in \texttt{RA\_D1\_Proofs.lean},
zero sorry tags).

Crucially, the LLC holds \emph{locally} on each causally connected
component.
Across causal severances---topological discontinuities that form
when the actualization flux reaches the bandwidth limit of the
local causal structure---conservation is not violated; it is
formally undefined, because there is no causal connection across
which to define it.
Event horizons are not information destroyers; they are boundaries
at which the graph has no outgoing edges.

% ── III. THE ENGINE OF BECOMING ─────────────────────────────────────────
\section{The Engine of Becoming}
\label{sec:engine}

The Engine of Becoming is a Covariant Stochastic Graph Rewriting
System stated at three levels of description.
The fundamental level (Level~1) is purely graph-theoretic.
The intermediate level (Level~2) assigns complex amplitudes to
causal histories via the quantum measure.
The emergent level (Level~3) recovers the smooth continuum in the
dense-graph limit.

The five steps of the algorithm are as follows.

\paragraph{Step 1 --- The Causal Boundary.}
The present frontier $\partialG$ is the set of actualized vertices
with no causal successors.
This is a causal notion, not a spatial one.
Different observers have different causal pasts and therefore different
$\partialG$, but the Level~1 growth rules are the same for all
observers.

\paragraph{Step 2 --- The Quantum Possibility Space.}
From $\partialG$, the valid candidate set $\Vvalid$ is determined
by the actualization criterion.
Amplitudes are assigned to all possible extensions via the quantum
measure~\cite{Sorkin1994}:
\begin{equation}
  \mu(\mathcal{A}) \;=\; \sum_{h,h'\in\mathcal{A}} D(h,h'),
  \label{eq:qm}
\end{equation}
where $D(h,h')$ is the decoherence functional over causal histories.
Virtual candidates carry zero metric weight (they are not vertices
of $V$) and do not source gravity.
This is the structural explanation for why the vacuum does not
gravitate: virtual (off-shell) processes satisfy $\Delta S = 0$
and add no vertices to $\DAG$.

\paragraph{Step 3 --- The Kinematic Snap.}
The actualization criterion filters candidates by requiring
$\Delta S(\rho\|\sigma_0) > 0$.
Because Lorentz boosts are unitary conjugations, and relative entropy
is invariant under unitary conjugation (proved in Lean:
\texttt{frame\_independence} in \texttt{RA\_AQFT\_Proofs\_v10.lean}),
this threshold is a Lorentz scalar.
A Rindler observer perceives the Minkowski vacuum as a KMS thermal
bath, but the relative entropy of that bath with respect to $\sigma_0$
is zero under modular flow (\texttt{rindler\_stationarity}, Lean-verified):
no actualization occurs.
Both observers agree.

\paragraph{Step 4 --- The Covariant Stochastic Update.}
A set of causally disconnected valid candidates
$\{v_1,\ldots,v_k\} \subseteq \Vvalid$ is actualized with probability
\begin{equation}
  P(\{v_1,\ldots,v_k\})
  \;=\; \left|\sum_\sigma \prod_i a\bigl(v_{\sigma(i)}
    \mid \mathrm{past}(v_{\sigma(i)})\bigr)\right|^2,
  \label{eq:update}
\end{equation}
where the sum is over all causal-order-consistent orderings $\sigma$.
By causal invariance~\cite{Sorkin1994}, this probability is the same
for every foliation: it depends only on the causal structure of the
antecedents.
The Born rule is not a postulate; it is a consequence of the quantum
measure on the causal graph.

\paragraph{Step 5 --- The Hydrodynamic Limit.}
After actualization, $V$ is updated and the effective metric is
recalculated from the actualization density $\lambda(x)$ via the
RA field equation:
\begin{equation}
  G_{\mu\nu} \;=\; 8\pi G\!\left(T_{\mu\nu}^{(A)}
    + \Theta_{\mu\nu}^{(\lambda)}\right),
  \label{eq:field_eq}
\end{equation}
where $\Theta_{\mu\nu}^{(\lambda)} = \xi[\nabla_\mu\nabla_\nu
(\ln\lambda) - g_{\mu\nu}\Box(\ln\lambda)]$.
This recalculation happens only \emph{after} actualization, never
during the quantum exploration.
Gravity is a post-actualization output of Step~5.
Superposed states never source the metric.

% ── IV. QM AND GR AS LIMITS ─────────────────────────────────────────────
\section{QM and GR as Limiting Approximations}
\label{sec:limits}

\subsection{Quantum mechanics as the slowly-varying limit}

When the actualization density $\lambda(x)$ varies slowly in space
and time---equivalently, when back-reaction of the causal graph onto
itself is negligible---the quantum measure of equation~\eqref{eq:qm}
reduces to the Feynman path integral:
\begin{equation}
  \langle x_f | e^{-iHt/\hbar} | x_i \rangle
  \;\underset{\lambda\to\mathrm{smooth}}{\longrightarrow}\;
  \int \mathcal{D}x\, e^{iS[x]/\hbar}.
  \label{eq:qm_limit}
\end{equation}
Standard quantum mechanics is the Engine of Becoming at Level~2
in the regime where the discrete causal graph is well-approximated
by a continuous Hilbert space evolution.

\subsection{General relativity as the dense-graph limit}

When the actualization density $\lambda \to \infty$, the discrete
graph converges to a smooth Lorentzian manifold and the RA field
equation~\eqref{eq:field_eq} reduces to:
\begin{equation}
  G_{\mu\nu} \;=\; 8\pi G\, T_{\mu\nu},
  \label{eq:gr_limit}
\end{equation}
the Einstein field equations.
The $\Theta_{\mu\nu}^{(\lambda)}$ correction vanishes when $\lambda$
is spatially uniform (dense homogeneous matter), recovering standard
GR in the interior of galaxies and on cosmological scales where
the density is approximately uniform.

\emph{Uniqueness of the EFE.}
The Einstein field equations with $\Lambda = 0$ are not merely
recovered but are the \emph{unique} macroscopic output of the
Engine of Becoming under three constraints: (i) the Local Ledger
Condition~\eqref{eq:LLC}, proved in Lean; (ii) the BDG action as
the unique discrete Lorentz-invariant action in $d=4$~\cite{Glaser2014};
and (iii) Lovelock's theorem~\cite{Lovelock1971}, which establishes
that the Einstein tensor is the unique rank-$(0,2)$ divergence-free
tensor constructed from the metric and its first two derivatives
in four spacetime dimensions.
No cosmological constant term arises because the vacuum does not
actualize (Step~2 of the Engine; proved in Paper~3~\cite{Sandeman2026P3}).
The combination of these three results closes the derivation without
free parameters.

\subsection{Why they appear incompatible}

QM assumes that states evolve continuously and unitarily (the
Level~2 quantum measure).
GR assumes that the spacetime metric is sourced by the stress-energy
tensor (the Level~3 dense-graph output).
These two approximations conflict in regimes where both matter
simultaneously---specifically, at Planck densities where quantum
fluctuations of the metric are non-negligible.
The apparent incompatibility is not a fundamental conflict between
two correct theories; it is the natural breakdown of two
approximations to the Engine of Becoming applied simultaneously
outside their domain of validity.

% ── V. EMERGENT COMPLEXITY AND THE CAUSAL FIREWALL ──────────────────────
\section{Emergent Complexity and the Causal Firewall}
\label{sec:complexity}

\subsection{Complexity as iterated coarse-graining}

Step~5 of the Engine of Becoming is recursive.
The hydrodynamic limit that produces the spacetime metric from
the actualization density $\lambda(x)$ is one instance of a
general coarse-graining operation.
Applying the same operation at successively larger scales produces
the full hierarchy of physical complexity.

Define the \emph{coarse-grained assembly depth} $\tilde{A}_{\rm RA}(x)$
as the mean causal depth of the DAG within a coarse-graining volume
$\mathcal{V}$ around $x$.
For a monoatomic gas, $\tilde{A}_{\rm RA} \sim 1$; for a covalent
molecule, $\tilde{A}_{\rm RA} \sim 3$--$5$; for a metabolically
active cell, $\tilde{A}_{\rm RA} \gg 10^3$.
The step from physics to chemistry to biology is the step from
low-depth to high-depth actualization networks.
This is not an emergentist metaphor; it is a structural consequence
of iterated Step~5 coarse-graining.

\subsection{Biological assembly and irreversibility}

Molecular assembly in biology corresponds directly to actualization
events.
Each JOIN operation in Assembly Theory~\cite{Cronin2023}---the
formation of a covalent bond that advances molecular construction---
corresponds to at least one irreversible actualization event:
an exothermic on-shell energy transfer that satisfies
$\Delta S > 0$ and cannot be reversed by subsequent unitary evolution.

This claim is not merely a formal identification.
B3LYP/6-311+G* density functional calculations across the five
primary bond types of biological assembly (C-C, C-N, C-O, S-S, P-O)
confirm that all are exothermic under standard biological conditions,
mandating irreversible energy dissipation at each assembly step.
The Local Ledger Condition~\eqref{eq:LLC} then governs the energy
bookkeeping of the assembly process, establishing a tight connection
between the physics of the DAG and the chemistry of life.

\subsection{The Causal Firewall}

The Engine of Becoming predicts a phase transition in the
actualization density that acts as a hard boundary separating
abiotic chemistry from stable biological complexity.

\begin{proposition}[Causal Firewall]
\label{prop:firewall}
At the critical actualization density $\mu = 1$ (the BDG self-dual
point, the QCD confinement--deconfinement transition), the
Erd\H{o}s-R\'{e}nyi random graph transitions from subcritical
($\mu < 1$: only small components, no spanning cluster) to
supercritical ($\mu > 1$: a giant connected component spans the
system).
This transition is the \emph{Causal Firewall}: below it, molecular
assemblies dissipate before reaching biological complexity; above
it, stable, self-reinforcing causal networks can persist.
\end{proposition}

The same critical density $\mu = 1$ governs four physically
distinct phenomena simultaneously: the QCD confinement--deconfinement
transition (particle physics); the onset of flat galactic rotation
curves (galactic dynamics via the $\nabla^2(\ln\lambda)$ term);
the fault-tolerance threshold of quantum error correction (quantum
information); and the Causal Firewall for biological complexity.
That a single dimensionless parameter controls the transition in
all four regimes is a structural prediction of the framework,
not a parameter choice.

\subsection{Biosignature prediction}

The Causal Firewall implies that substrate-independent physical
constraints govern the emergence of life.
Specifically: any planetary environment in which biological
complexity can arise must sustain actualization densities above
the Causal Firewall threshold.
This is independent of the chemical substrate (carbon-based,
silicon-based, or otherwise).

The falsifiable prediction: biosignature searches should correlate
the presence of biological complexity with thermodynamic environments
that sustain $\mu \geq 1$.
Environments that are too cold, too rarefied, or too far from
chemical disequilibrium cannot sustain biological assembly because
they fall below the Causal Firewall threshold.
This is a physically derived constraint on the distribution of life
in the universe, testable with current and near-future
astrobiology surveys.

% ── VI. TIME AND THE ARROW ───────────────────────────────────────────────
\section{Time, the Arrow, and Ontic Randomness}
\label{sec:time}

\subsection{Time as the growing graph}

In RA, time is not a dimension through which events are located.
Time \emph{is} the process of the causal DAG growing.
Proper time along a worldline is the count of actualization events
along that worldline:
\begin{equation}
  \tau_{\rm proper} \;=\; \frac{N_{\rm act}}{\lambda_0},
  \label{eq:proper_time}
\end{equation}
where $N_{\rm act}$ is the number of actualization events and
$\lambda_0$ is the reference actualization rate in flat space.
Relativistic time dilation is the statement that different worldlines
accumulate actualization events at different rates: a clock in
motion writes fewer vertices per coordinate second than a stationary
clock, because its causal graph is more sparsely connected to the
environment~\cite{Sandeman2026P3}.

This identification is not a metaphor.
The Lorentz transformation acts on the DAG as a re-labelling of
vertices; proper time is invariant under this re-labelling because
it counts physical events, not coordinate intervals.

\subsection{The arrow of time}

Thermodynamic accounts of the arrow of time require explaining why
entropy increases from past to future given time-symmetric
microscopic laws.
In RA, the arrow of time is not derived from entropy; it is
\emph{encoded in the definition of an actualization event}.
An edge $(u, v) \in E$ is directed: $u$ is in the causal past of
$v$, and this relation cannot be reversed.
There is no sense in which the DAG can shrink.
The past is permanent; the future is open.

This dissolves the problem.
One cannot ask why the DAG grows in the direction it does:
directedness is constitutive of what an edge is.
The arrow of time is not a feature to be explained by the framework;
it is the framework's most primitive axiom made explicit.

\subsection{Ontic randomness}

The randomness of quantum measurement is ontically irreducible in RA.
In a two-state superposition $\alpha|x_1\rangle + \beta|x_2\rangle$,
the Born rule probabilities $|\alpha|^2$, $|\beta|^2$ are the
distribution over possible next vertices of the DAG.
Which vertex is added at Step~4 is not determined by any prior
fact in $V$: it is a genuinely new addition to the causal order.
The actualization event at $x_0$ is the first link of a new causal
chain, not the conclusion of an old one.

This is not ignorance of a hidden variable.
The quantum amplitudes governing the next vertex are not waiting to
be discovered; they are underdetermined until the event occurs.
Classical chaos is different: given perfect initial conditions,
the trajectory is fixed; only our ignorance produces apparent
randomness.
In RA, there are no initial conditions that would determine the
outcome if we knew them.
The randomness is in the world, not in our description of it.

% ── VI. PHILOSOPHICAL POSITIONING ──────────────────────────────────────
\section{Philosophical Positioning}
\label{sec:philosophy}

\subsection{Against the block universe}

The block universe (eternalism) holds that past, present, and future
events are equally real, and that temporal becoming is an illusion
produced by our position within the four-dimensional manifold.
In RA, the block universe is literally false: the DAG $\DAG$ is
not a complete static structure but a growing one.
Future vertices do not exist; they are not yet actual.
The distinction between the actual and the possible is metaphysically
sharp, not perspectival.
This is not a claim about human experience of time but about the
physical structure of the world.

\subsection{Against Copenhagen}

The Copenhagen interpretation grounds measurement in the classical/quantum
distinction without specifying what draws the boundary.
RA provides the criterion: the boundary lies where $\Delta S > 0$.
There is no Heisenberg cut in RA; there is only a gradient in the
actualization rate $\lambda(x)$.
Large, warm objects are classical not because they are fundamentally
different from quantum systems but because their internal actualization
rate is far above the threshold at every observable scale.

\subsection{Against Everett}

Everettian QM holds that all branches of the wave function are real
and that apparent collapse is a perspectival artifact.
RA holds that exactly one branch becomes actual at each Step~4;
the others were genuinely possible but are now closed.
The difference is not merely semantic.
Everett predicts no sharp reversibility threshold (all branches
remain in principle recoverable); RA predicts a specific measurable
threshold $\Delta S^*$ above which recovery is impossible.
Everett predicts no Lieb-Robinson propagation delay for apparent
collapse; RA predicts a finite propagation speed~\cite{Sandeman2026P2}.
These are discriminating empirical differences.

\subsection{Against pilot wave theory}

Bohmian mechanics accounts for quantum randomness by postulating
hidden particle positions guided by the wave function.
RA posits no hidden variables.
The randomness at Step~4 is irreducibly ontic, not epistemic:
there is no fact of the matter about which vertex will be added
before it is added.
Bohmian mechanics is also in tension with quantum field theory,
where particle number is not conserved; RA operates at the level
of field interactions and is free of this constraint.

\subsection{Against relational QM}

Rovelli's relational QM~\cite{Rovelli1996} grounds quantum
states relative to observers.
RA's actualization criterion is observer-independent: the relative
entropy $\Delta S(\rho\|\sigma_0)$ is a property of the field state
with respect to the physical vacuum, not with respect to any observer.
The Rindler resolution (Section~\ref{sec:engine}) makes this explicit:
both the inertial and Rindler observers agree on $\Delta S = 0$ for
the vacuum, regardless of their different experiences.

% ── VII. PREDICTIONS ────────────────────────────────────────────────────
\section{Falsifiable Predictions}
\label{sec:predictions}

Table~\ref{tab:all_predictions} assembles the complete set of
falsifiable quantitative predictions of the Relational Actualism
programme.
All predictions are structural consequences of the single ontological
commitment of Section~\ref{sec:primitive}; none are post-hoc fits.

\begin{table*}[t]
\centering
\caption{Complete falsifiable prediction set of the Relational Actualism
  programme. All predictions follow structurally from the single commitment
  that actualization events are primitive. Paper numbers refer to the
  companion papers~\cite{Sandeman2026P1,Sandeman2026P2,Sandeman2026P3}.
  $\checkmark$ = confirmed by existing data; $\bigcirc$ = testable with
  current or near-future experiments; $\square$ = requires future
  infrastructure.}
\label{tab:all_predictions}
\small
\begin{tabular}{p{4.5cm}p{3.5cm}lp{3cm}}
\toprule
Prediction & Value & Status & Source / Test \\
\midrule
\multicolumn{4}{l}{\textit{Paper 1: BDG causal geometry}} \\[2pt]
$\alpha_{\rm EM}^{-1}$ (leading order) & $144 - 7 = 137.000$ & $\checkmark$ 0.026\% & PDG 2024 \\
$\alpha_{\rm EM}^{-1}$ (exact polynomial) & 136.845 & $\checkmark$ 0.14\% & PDG 2024 \\
$\alpha_s(m_Z)$ & $1/\sqrt{72} = 0.11785$ & $\checkmark$ 0.13\% & PDG 2024 \\
$m_p / \Lambda_{\rm QCD}$ & $\sqrt{c_4} = 2\sqrt{2}$ & $\checkmark$ 0.08\% & PDG 2024 \\
SM particle content & 5 topology classes & $\checkmark$ exact & Lean-verified \\
No proton decay & $\tau_p = \infty$ & $\checkmark$ (lower bounds) & Super-K, Hyper-K \\
No 4th SM generation & 3 generations only & $\checkmark$ & LEP, LHC \\
\midrule
\multicolumn{4}{l}{\textit{Paper 2: Wave function collapse}} \\[2pt]
LR actualization propagation delay & $\Delta t \geq |r_A - r_B|/v_{\rm LR}$ & $\bigcirc$ & Optical lattices, SC qubits \\
Env-dependent transition duration & $\tau_{\rm act} \sim \hbar/\Delta E_{\rm env}$ & $\bigcirc$ & SC circuit tuning \\
Sharp reversibility threshold & $\Delta S^*$ measurable & $\bigcirc$ & Quantum eraser+threshold \\
Rindler stationarity & $\Delta S = 0$ under modular flow & $\checkmark$ Lean-verified & Unruh consistency \\
\midrule
\multicolumn{4}{l}{\textit{Paper 3: Cosmology}} \\[2pt]
$H_0$ in Eridanus supervoid & $\approx 75.9\;\mathrm{km\,s^{-1}\,Mpc^{-1}}$ & $\bigcirc$ & Void surveys \\
$H_0$ anticorrelates with $\bar\rho_{\rm LOS}$ & line-of-sight dependence & $\bigcirc$ & DESI, Euclid \\
No WIMP gravitational sourcing & $\leq 10^{-25}$ of baryons & $\checkmark$ null results & LZ, PandaX-4T \\
MW-analog Einstein radius & $\approx 0.47$ arcsec & $\bigcirc$ & Strong lensing surveys \\
Structural $\Lambda = 0$ & no vacuum contribution & $\bigcirc$ & CMB low-$\ell$ power \\
\midrule
\multicolumn{4}{l}{\textit{This paper (P4)}} \\[2pt]
BMV null result & no gravity-mediated entanglement & $\bigcirc$ & BMV experiment \\
Kinematic Coherence Bound & $N_{\rm max} = \eta\cdot p_{\rm th}$ & $\bigcirc$ & Fault-tolerant QC arrays \\
No Boltzmann Brains & $P_{\rm BB} = 0$ exactly & $\checkmark$ (structural) & Vacuum definition \\
\bottomrule
\end{tabular}
\end{table*}

\subsection{The BMV prediction}

The Bose-Marletto-Vedral (BMV) experiment~\cite{Bose2017,Marletto2017}
proposes to measure gravitationally mediated entanglement between
two masses in superposition.
RA predicts a strict null result: because the metric is updated only
at Step~5, after actualization, and quantum superpositions never
source the metric, there can be no gravitationally mediated
entanglement between systems that have not already undergone
actualization (Section~\ref{sec:engine}).
A null result would be the first direct experimental support for the
framework's core architecture.
A positive result would falsify the Engine of Becoming as stated.
Both outcomes are scientifically valuable.

\subsection{The Kinematic Coherence Bound}

The actualization criterion imposes a maximum scale on fault-tolerant
quantum computation.
An array of $N$ physical qubits with error rate $p$ per gate can
maintain logical coherence only up to
\begin{equation}
  N_{\rm max} \;\leq\; \eta \cdot p_{\rm th},
  \label{eq:KCB}
\end{equation}
where $p_{\rm th}$ is the fault-tolerance threshold and $\eta$
is a constant determined by the actualization geometry.
This bound is absent from standard QEC theory, which imposes no
upper limit on array size given sufficiently small error rates.
Testing this bound at large qubit array scales provides a direct
probe of the actualization threshold.

\subsection{Boltzmann Brains}

The Boltzmann Brain problem asks whether the universe is dominated
by isolated thermodynamic fluctuations (``brains'') in an otherwise
empty de Sitter vacuum.
In RA, spontaneous fluctuations from the vacuum satisfy
$\Delta S(\rho\|\sigma_0) = 0$ (the vacuum state is the reference
$\sigma_0$ itself), and therefore produce no actualization events.
Boltzmann Brains are categorically impossible: they require
self-organizing complexity to emerge from a state that, by definition,
contains no actualization events.
This is a structural prohibition, not a probabilistic suppression.

% ── VIII. HARD WALLS ────────────────────────────────────────────────────
\section{Open Problems and Hard Walls}
\label{sec:hardwalls}

The following mathematical problems are precisely identified gaps
in the present formulation.
They do not undermine the predictions in Section~\ref{sec:predictions},
which hold in the regimes where the existing derivations apply.

\paragraph{Hard Wall 1: Bianchi compatibility.}
The RA field equation~\eqref{eq:field_eq} must satisfy the contracted
Bianchi identity $\nabla_\mu G^{\mu\nu} \equiv 0$, requiring
$\nabla_\mu\Theta^{\mu\nu}_{(\lambda)} = 0$.
This is satisfied in the weak-field static limit; the full covariant
proof requires deriving the coupling $\xi$ from the microscopic graph
dynamics.
Collaboration target: Algebraic QFT expertise in Tomita-Takesaki
modular theory.

\paragraph{Hard Wall 2: Amplitude locality.}
The causal invariance theorem at Step~4 is proved in the perturbative
regime given the axiom of amplitude locality.
The intrinsic discrete proof---showing that local amplitudes satisfy
causal independence as a theorem of the discrete graph dynamics
rather than an axiom---is the primary open problem of the Lean
verification programme.
This is equivalent to proving Bell causality for the graph rewriting
system in the Rideout-Sorkin~\cite{RideoutSorkin2000} sense.
Collaboration target: Causal set theory experts (Dowker, Henson).

\paragraph{Hard Wall 3: Continuum limit.}
The discrete LLC~\eqref{eq:LLC} is Lean-verified.
That the continuum limit satisfies the contracted Bianchi identity
$\nabla_\mu G^{\mu\nu} \equiv 0$ without breaking Lorentz invariance
requires demonstrating convergence of the discrete conservation
law to the differential Noether condition.
Collaboration target: Experts in causal dynamical triangulations
(Loll, Radboud) and asymptotic safety.

\paragraph{Hard Wall 4: Covariant $\Box(\ln\lambda)$.}
The cosmological predictions of Paper~3 use the weak-field
$\nabla^2(\ln\lambda)$ correction.
The fully covariant $g^{\mu\nu}\nabla_\mu\nabla_\nu(\ln\lambda)$
is required for strong-field and early-universe applications.
Demonstrating that this satisfies Bianchi compatibility and
reduces to the Newtonian limit is Hard Wall~2 of the cosmology paper.

% ── IX. CONCLUSION ──────────────────────────────────────────────────────
\section{Conclusion}
\label{sec:conclusion}

Relational Actualism begins from a single claim: on-shell actualization
events are the primitive constituents of physical reality.
From this commitment alone, with no additional fields or free
parameters, the programme derives: both Standard Model coupling
constants and the complete particle spectrum (Paper~1); an
observer-independent criterion for wave function collapse with three
testable predictions (Paper~2); the structural zero of the
cosmological constant, the Hubble tension, and dark matter
phenomenology (Paper~3); and the framework presented here,
which shows that QM and GR are approximations to the same algorithm.

The arc of the argument is conceptually clean.
The measurement problem is dissolved because collapse is identified
with a physical event (irreversible entropy production), not an
observer's intervention.
The cosmological constant problem is dissolved because the vacuum
does not gravitate (off-shell processes add no vertices to the DAG).
The Hubble tension has a specific geometric mechanism (expansion is
inversely proportional to actualization density).
The dark matter phenomenology is reproduced by the topological
tension at density gradients in the causal graph.

In each case, the resolution is not ad hoc.
It is the same ontological commitment applied consistently.

The universe may not always know what it is going to do, but it
always knows what it has done.
That permanence of the actual, and openness of the possible, is the
simplest---and perhaps the only---account of physical reality that
is simultaneously consistent with quantum mechanics, general
relativity, and the irreversible macroscopic world we inhabit.

% ── ACKNOWLEDGEMENTS ────────────────────────────────────────────────────
\section*{Acknowledgements}
The author acknowledges the use of Anthropic's Claude and Google's Gemini
as AI research partners during the preparation of this manuscript.
The author is solely responsible for all scientific content and conclusions.

% ── REFERENCES ──────────────────────────────────────────────────────────
\begin{thebibliography}{99}

\bibitem{Sorkin1994}
R.~D.~Sorkin,
Quantum mechanics as quantum measure theory,
Mod.\ Phys.\ Lett.\ A \textbf{9}, 3119 (1994).

\bibitem{RideoutSorkin2000}
D.~P.~Rideout and R.~D.~Sorkin,
Classical sequential growth dynamics for causal sets,
Phys.\ Rev.\ D \textbf{61}, 024002 (2000).

\bibitem{Rovelli1996}
C.~Rovelli,
Relational quantum mechanics,
Int.\ J.\ Theor.\ Phys.\ \textbf{35}, 1637 (1996).

\bibitem{Bose2017}
S.~Bose et al.,
Spin entanglement witness for quantum gravity,
Phys.\ Rev.\ Lett.\ \textbf{119}, 240401 (2017).

\bibitem{Marletto2017}
C.~Marletto and V.~Vedral,
Gravitationally induced entanglement between two massive particles
is sufficient evidence of quantum effects in gravity,
Phys.\ Rev.\ Lett.\ \textbf{119}, 240402 (2017).

\bibitem{Sandeman2026P1}
J.~F.~Sandeman,
The Fine Structure Constant and Strong Coupling from
Four-Dimensional Causal Geometry,
Zenodo doi:10.5281/zenodo.19324790 (2026).

\bibitem{Sandeman2026P2}
J.~F.~Sandeman,
Wave Function Collapse as Irreversible Entropy Production:
An Observer-Independent Criterion with Three Experimental Predictions,
in preparation (2026).

\bibitem{Sandeman2026P3}
J.~F.~Sandeman,
Structural Zero Cosmological Constant and the Hubble Tension from
Causal Actualization Density,
in preparation (2026).



\bibitem{Glaser2014}
L.~Glaser,
A closed form expression for the causal set d'Alembertian,
Class.\ Quantum Grav.\ \textbf{31}, 095007 (2014).

\bibitem{Lovelock1971}
D.~Lovelock,
The Einstein tensor and its generalizations,
J.\ Math.\ Phys.\ \textbf{12}, 498 (1971).

\bibitem{Cronin2023}
S.~M.~Walker et al.,
Identifying molecules as biosignatures with assembly theory
and mass spectrometry,
Nature \textbf{616}, 01 (2023).

\end{thebibliography}

\end{document}
